定义有向边方向上的子树哈希
\[
  h(u,p)=1+\sum_{v\in g[u],\,v\ne p}
  \operatorname{mix}\bigl(h(v,u)\bigr)
  \pmod{2^{64}}.
\]
代码使用 \texttt{uint64\_t} 自然溢出，不需要手动取模。
求和使儿子顺序无关，同时保留相同儿子出现的次数。

代码块直接放在 \texttt{g} 建好之后。
\texttt{dfs(u,p)} 返回以 $u$ 为根且不经过 $p$ 的子树哈希；整棵树以
\texttt{root} 为根的哈希就是 \texttt{dfs(root,0)}。

每个儿子哈希先异或随机盐，经过三步 xorshift 后再异或一次盐。
同一次程序运行中的盐相同，因此可以直接比较多棵树；重新运行后哈希值通常会改变。
该方法单次 DFS 为 $O(n)$，属于概率哈希，理论上可能碰撞。
